
\documentclass[10pt]{article}
\usepackage[utf8]{inputenc}
\usepackage[a4paper, left=2.54cm, right=2.54cm, top=3cm, bottom=3cm]{geometry} % Márgenes personalizados a 2cm
\usepackage[spanish]{babel} % Para fechas y formato en español
\usepackage{multicol}
\usepackage{enumitem}
\usepackage{fancyhdr}
\usepackage{titlesec}

% Configuración de página
\pagestyle{fancy}
\fancyhf{} % Limpia encabezados y pies de página
\fancyhead[L]{Academia Prepolicial - Números Racionales}
\fancyhead[R]{\today}
\fancyfoot[C]{\thepage} % Número de página centrado en el pie
\renewcommand{\headrulewidth}{0.8pt}

% Configuración de títulos y espaciado
\titleformat{\section}[block]{\large\bfseries\centering}{\thesection}{1em}{}
\setlength{\columnsep}{1cm}
\setlength{\parskip}{0.8em}
\setlength{\parindent}{0em}


\begin{document}

	\begin{center}
		\Large\bfseries Números Racionales
	\end{center}

\begin{multicols}{2}
	\section*{Ejercicios}
		\begin{enumerate}[label=\textbf{\arabic*.}, itemsep=0.4cm]
			\item Calcula la diferencia entre los términos de una fracción equivalente a \(\frac{2}{5}\), sabiendo que la suma de dichos términos es 28.

			\item Calcula la suma de los términos de la fracción irreductible equivalente a \(\frac{420}{924}\).

			\item Si al numerador de una fracción irreductible se le suma 1 y al denominador se le suma 2, resulta ser equivalente a la fracción original. La fracción original es:

			\item \textbf{PUCP}
			Calcula la suma de todos los valores de \(a\) sabiendo que la fracción \(\frac{a}{12}\) es propia e irreductible.

			\item Calcula la suma de todos los valores de \(a\) sabiendo que la fracción \(\frac{a}{18}\) es propia e irreductible.

			\item Calcula el número de fracciones equivalentes a \(\frac{4}{5}\), si su numerador está comprendido entre 25 y 40 y su denominador entre 38 y 53.

			\item Calcula la suma de cifras del numerador de una fracción cuyo denominador es 84, sabiendo además que dicha fracción es mayor que \(\frac{1}{7}\) pero menor que \(\frac{1}{6}\).
			\textbf{UNMSM}

			\item ¿Cuál es la diferencia de los términos de la fracción equivalente a \(\frac{3}{5}\) cuya suma de términos es 9696?

			\item ¿Cuál es la diferencia de los términos de la fracción equivalente a \(\frac{7}{8}\) cuya suma de términos es 555?

			\item Determina una fracción cuya suma de términos es 25 y cuando se le suma 6 unidades al numerador y 9 al denominador, se obtiene una fracción equivalente a \(\frac{3}{5}\). Dar como respuesta la diferencia de los 2 términos de la fracción.

			\item Si \(0,\overline{1a} = \frac{m}{11}\), calcula: \(a + m\).


			\item El denominador excede al numerador de una fracción en la unidad. Si al denominador se le agrega 4 unidades, el resultado es 2 unidades menos que el triple de la fracción original. ¿Cuál es el numerador de la fracción original?

			\item Encuentra una fracción que si a sus 2 términos se les suma el denominador y al resultado se le resta la fracción, se obtenga la misma fracción.

			\item Si a los dos términos de una fracción irreductible se les suma el cuádruple del denominador y al resultado se le resta la fracción, resulta la misma fracción. ¿Cuánto suman los términos de la fracción original?


			% Para añadir más ejercicios, simplemente copia y pega el siguiente bloque:
			% \item [Enunciado del ejercicio]
		\end{enumerate}

\end{multicols}

\end{document}